% titel: Blatt 10, A7: Na2HPO4/KH2PO4-Puffer, Menge KH2PO4
% Teil II: der durchgerechnete Weg.

\subsection*{Gegeben und gesucht}
Puffer aus Na$_2$HPO$_4$ (Base, HPO$_4^{2-}$) und KH$_2$PO$_4$ (Säure,
H$_2$PO$_4^-$), $pK_s=7.21$, Ziel-pH $=7.0$. 50 mL einer 0.5 M
Na$_2$HPO$_4$-Lösung. Gesucht: Masse KH$_2$PO$_4$.

\subsection*{Ansatz}
$pH=pK_s+\lg\frac{[\text{Base}]}{[\text{Säure}]}$;
$M(\mathrm{KH_2PO_4})\approx136.09\ g\,mol^{-1}$.

\subsection*{Rechnung}
Stoffmenge Base: $n_{\text{Base}}=0.5\cdot0.050=0.025$ mol.
\begin{align*}
  7.0 &= 7.21 + \lg\frac{[\text{Base}]}{[\text{Säure}]}
    \Rightarrow \lg\frac{[\text{Base}]}{[\text{Säure}]} = -0.21,\\
  \frac{[\text{Base}]}{[\text{Säure}]} &= 10^{-0.21} \approx 0.617 .
\end{align*}
Da gleiches Volumen, gilt das Verhältnis auch für die Mengen:
\[
  n_{\text{Säure}} = \frac{n_{\text{Base}}}{0.617}
    = \frac{0.025}{0.617} \approx 0.0405\ mol .
\]
\[
  m(\mathrm{KH_2PO_4}) = n\,M = 0.0405\cdot136.09 \approx 5.5\ g .
\]

\begin{ergebnis}
Man benötigt ca.\ $5.5$ g Kaliumhydrogenphosphat (KH$_2$PO$_4$) zur
Na$_2$HPO$_4$-Lösung, um pH 7.0 zu erreichen.
\end{ergebnis}

\begin{kontrolle}
$pH<pK_s$ $\Rightarrow$ im Puffer überwiegt die Säure; $n_{\text{Säure}}>
n_{\text{Base}}$ (0.0405 > 0.025) bestätigt das. Probe:
$7.21+\lg\frac{0.025}{0.0405} = 7.21+\lg 0.617 =7.21-0.21=7.0$. Passt.
\end{kontrolle}

\begin{merksatz}
Henderson-Hasselbalch: $pH=pK_s+\lg\frac{[\text{Base}]}{[\text{Säure}]}$.
Aus $pH<pK_s$ folgt Säure-überschuss; Mengenverhältnis $=10^{pH-pK_s}$
verwenden.
\end{merksatz}
